\documentclass[11pt]{article}
\usepackage{amsmath,amssymb,amsthm}
\usepackage[T1]{fontenc}
\usepackage{amsmath,amssymb,amsthm}
\usepackage[margin=1in]{geometry}
\usepackage{microtype}

\newtheorem{lemma}{Lemma}
\newtheorem{proposition}[lemma]{Proposition}
\theoremstyle{definition}
\newtheorem{definition}[lemma]{Definition}
\theoremstyle{remark}
\newtheorem*{remark}{Remark}

\newcommand{\R}{\mathbb{R}}
\newcommand{\htop}{h_{\mathrm{top}}}
\newcommand{\Dprime}{\mathcal{D}'}

\newcommand{\fptitle}{Review of Problem 8 -- Astra}
\newcommand{\reviewtag}{Reviewer 1}
\usepackage{fp-review}

\begin{document}
\maketitle

\section{Review}

\paragraph{Correctness.}
\begin{itemize}
	\item  1: essentially flawless mathematical argument
	\item  0.5: minor errors that are easily corrected
	\item  0: blank solution, or acknowledged partial result
	\item  $-0.5$: major errors
	\item  $-1$: proposed proof makes no progress towards the result.
\end{itemize}
Acknowledges that it doesn't solve problem.

\noindent Score: 0

\paragraph{Exposition.}
\begin{itemize}
	\item  1: well-written and clear solution.
	\item  0: argument that is possible to follow with serious effort
	\item  $-1$: very poor exposition
\end{itemize}

\noindent Score:0

\paragraph{Attributions.}
\begin{itemize}
	\item  1: solutions contains fully appropriate attributions
		to related work and theorems being used.
	\item  0: moderate issues with missing attributions or hallucinated references
	\item  $-1$: blatant plagiarism or hallucinations
\end{itemize}

\noindent Score: 0
\review{The implication from flow to diffeomorphism has been known for ever though suspension, it is a simple exercise and it is already mentioned in a paper of the reviewer with coauthor, so there is not much value in this step. Even the other direction would represent some interest, though quite unrelated with the main problem. Exposition is not bad, easy reading but mostly because mathematics is trivial.}
\newpage
\section{Submission with inline comments}


\begin{abstract}
A complete proof of the proposed assertion is not known in the stated
generality. In fact, the assertion would settle the open problem of
vanishing Lyapunov exponents for cohomologically rigid diffeomorphisms.
We give a complete proof of this reduction, explaining why the usual
structural results about globally hypoelliptic vector fields do not
supply the requested conclusion. 

\end{abstract}

\section{Status of the question}

The hypothesis in the question is called \emph{global hypoellipticity}.
The established theory relates this property to cohomology freeness and
gives unique ergodicity with a smooth invariant probability density;
see Forni~\cite[Proposition~2.4 and Theorem~3.10]{Forni}.
It does not, however, supply the proposed general entropy conclusion.

The relevant unresolved issue is explicitly recorded by
Karaliolios~\cite[Introduction, p.~6]{Karaliolios}: it is unknown whether
cohomological rigidity of a diffeomorphism implies that all its Lyapunov
exponents vanish. The suspension construction below shows that a proof
of the assertion in the question would imply precisely this vanishing
result.

Thus the following is a rigorous reduction to an open problem, rather
than a purported proof of the requested assertion.

\section{A sufficient condition for global hypoellipticity}

We use smooth positive densities, so that orientability is unnecessary.
If $\mu$ is such a density, scalar distributions are identified with
continuous linear functionals on $C^\infty(M)$ through the pairing
\[
 \langle u,\phi\rangle_\mu
 =
 \langle u,\phi\mu\rangle.
\]
For a vector field preserving $\mu$, the distributional derivative
satisfies
\[
 \langle Xu,\phi\rangle_\mu
 =
 -\langle u,X\phi\rangle_\mu.
\]

\begin{lemma}\label{lem:cf-gh}
Let $Y$ be a smooth vector field on a compact manifold $P$, preserving
a smooth positive probability density $\mu$. Suppose that for every
$f\in C^\infty(P)$ there is $v\in C^\infty(P)$ such that
\begin{equation}\label{eq:flow-cf}
 Yv=f-\int_P f\,d\mu.
\end{equation}
Then $Y$ satisfies the global hypoellipticity hypothesis of the question.
\end{lemma}

\begin{proof}
Suppose that $Yu=f$ for $u\in\Dprime(P)$ and $f\in C^\infty(P)$.
Invariance of $\mu$ gives
\[
 \int_P f\,d\mu
 =
 \langle Yu,1\rangle_\mu
 =
 -\langle u,Y1\rangle_\mu
 =
 0.
\]
Consequently, \eqref{eq:flow-cf} supplies a smooth function $v$ with
$Yv=f$. Set $w=u-v$, so that $Yw=0$.

For any $\phi\in C^\infty(P)$, choose $a\in C^\infty(P)$ satisfying
\[
 Ya=\phi-\int_P\phi\,d\mu.
\]
Then
\[
 \langle w,\phi\rangle_\mu
 =
 \left(\int_P\phi\,d\mu\right)\langle w,1\rangle_\mu
 +\langle w,Ya\rangle_\mu
 =
 \left(\int_P\phi\,d\mu\right)\langle w,1\rangle_\mu.
\]
Thus $w$ is a constant distribution. It follows that $u=v+w$ is smooth.
\end{proof}

\section{Suspending a cohomologically rigid diffeomorphism}

\begin{definition}
Let $F:N\to N$ be a smooth diffeomorphism of a compact connected
manifold, preserving a smooth positive probability density $\nu$.
We call $F$ \emph{cohomologically rigid} if, for every
$a\in C^\infty(N)$ with $\int_N a\,d\nu=0$, there is
$v\in C^\infty(N)$ satisfying
\begin{equation}\label{eq:map-cf}
 v\circ F-v=a.
\end{equation}
\end{definition}

\begin{lemma}\label{lem:suspension}
The constant-roof suspension of a cohomologically rigid diffeomorphism
has a globally hypoelliptic generating vector field.
\end{lemma}

\begin{proof}
Let
\[
 P_F=(N\times\R)/{\sim},
 \qquad
 (x,s+1)\sim(Fx,s).
\]
The equivalence relation is generated by the free, properly
discontinuous smooth $\mathbb Z$-action
\[
 k\cdot(x,s)=(F^k x,s-k).
\]
Hence $P_F$ is a smooth manifold. It is compact because the image of
$N\times[0,1]$ is all of $P_F$, and it is connected because $N$ is
connected.

Translation in the second coordinate descends to the smooth flow
\[
 G_t[x,s]=[x,s+t],
\]
whose generating vector field $Y$ is induced by $\partial_s$.
Since $F$ preserves $\nu$, the product density $d\nu\,ds$ descends to
a smooth positive $G_t$-invariant probability density $\mu$ on $P_F$.

We verify \eqref{eq:flow-cf}. Given $f\in C^\infty(P_F)$, denote its
lift to $N\times\R$ by $\widetilde f$. Thus
\[
 \widetilde f(x,s+1)=\widetilde f(Fx,s).
\]
Set
\[
 c=\int_{P_F}f\,d\mu,
 \qquad
 a(x)=\int_0^1\bigl(\widetilde f(x,r)-c\bigr)\,dr.
\]
By the definition of the suspension measure,
$\int_N a\,d\nu=0$. Cohomological rigidity supplies
$v\in C^\infty(N)$ with $v\circ F-v=a$.

Define, for all $s\in\R$,
\[
 \widetilde u(x,s)
 =
 v(x)+\int_0^s\bigl(\widetilde f(x,r)-c\bigr)\,dr.
\]
This is smooth. Moreover,
\[
 \frac{\partial}{\partial s}
 \bigl(\widetilde u(x,s+1)-\widetilde u(Fx,s)\bigr)=0,
\]
and at $s=0$ the expression in parentheses equals
\[
 v(x)+a(x)-v(Fx)=0.
\]
Therefore
\[
 \widetilde u(x,s+1)=\widetilde u(Fx,s).
\]
It follows that $\widetilde u$ descends to a smooth function $u$ on
$P_F$, and
\[
 Yu=f-c=f-\int_{P_F}f\,d\mu.
\]
Lemma~\ref{lem:cf-gh} now proves global hypoellipticity.
\end{proof}

\section{What the requested assertion would imply}

\begin{proposition}\label{prop:reduction}
If the assertion in the question holds for all compact connected smooth
manifolds, then every cohomologically rigid smooth diffeomorphism
preserving a smooth positive probability density has zero topological
entropy and all its Lyapunov exponents vanish almost everywhere.
\end{proposition}

\begin{proof}
Let $F:N\to N$ be such a diffeomorphism and form the suspension from
Lemma~\ref{lem:suspension}. Its generating vector field satisfies
exactly the hypothesis of the question.

The embedded section
\[
 S=\{[x,0]:x\in N\}\subset P_F
\]
is invariant under the time-one map, since
\[
 G_1[x,0]=[x,1]=[Fx,0].
\]
Thus $G_1|_S$ is conjugate to $F$. Topological entropy cannot increase
upon restriction to a compact invariant subset: this follows directly
from the definition using separated sets, since every separated set
in the subset is also separated in the ambient space. Hence
\begin{equation}\label{eq:entropy-restriction}
 \htop(F)\leq\htop(G_1).
\end{equation}
The entropy of a flow is, by convention, that of its time-one map.
Consequently, the requested assertion would give
$\htop(G_1)=0$, and \eqref{eq:entropy-restriction} would imply
$\htop(F)=0$.

Metric entropy is bounded above by topological entropy
\cite{Goodwyn}, so
\[
 h_\nu(F)=0.
\]
Let $\lambda_1(x),\ldots,\lambda_d(x)$ denote the Lyapunov exponents,
listed with multiplicity. Since $F$ is smooth and $\nu$ is a smooth
invariant probability density, Pesin's entropy formula applies
\cite{Pesin}:
\[
 h_\nu(F)
 =
 \int_N\sum_{j=1}^d\max\{\lambda_j(x),0\}\,d\nu(x).
\]
The integrand is nonnegative. Its integral being zero, every Lyapunov
exponent is nonpositive at $\nu$-almost every point.

Finally, choose a Riemannian volume density $m$ and write
$d\nu=\rho\,dm$, where $\rho$ is smooth and everywhere positive.
Invariance of $\nu$ gives
\[
 \bigl|\det D_xF^n\bigr|
 =
 \frac{\rho(x)}{\rho(F^n x)}.
\]
Because $\rho$ is bounded above and bounded away from zero,
\[
 \sum_{j=1}^d\lambda_j(x)
 =
 \lim_{n\to\infty}\frac1n
       \log\bigl|\det D_xF^n\bigr|
 =
 0
\]
at every Lyapunov-regular point. A finite collection of nonpositive
numbers with sum zero consists entirely of zeros. Thus all Lyapunov
exponents vanish $\nu$-almost everywhere.
\end{proof}

\begin{remark}[The essential gap in a tempting entropy argument]
Preservation of a smooth volume gives
$\sum_j\lambda_j=0$, not
$\sum_j\max\{\lambda_j,0\}=0$.
The latter is the quantity occurring in Pesin's entropy formula.
Consequently, replacing entropy by the integral of the logarithm of
the full Jacobian determinant would be incorrect.
\end{remark}

\section{Conclusion}

Proposition~\ref{prop:reduction} shows that the question includes the
Lyapunov-exponent vanishing problem explicitly identified as open in
\cite{Karaliolios}. The established global-hypoellipticity results
provide cohomological rigidity and a smooth invariant measure, but
not the additional vanishing needed for entropy. Therefore a complete
proof of the proposed general assertion cannot presently be supplied
as an established mathematical result; doing so would require resolving
this open problem.

\begin{thebibliography}{9}

\bibitem{Forni}
G.~Forni,
\emph{On the Greenfield--Wallach and Katok conjectures in dimension three},
in \emph{Geometric and Probabilistic Structures in Dynamics},
Contemporary Mathematics \textbf{469},
American Mathematical Society, 2008, pp.~197--213.
See also arXiv:0706.3981, especially Proposition~2.4 and Theorem~3.10.

\bibitem{Goodwyn}
L.~W.~Goodwyn,
\emph{Topological entropy bounds measure-theoretic entropy},
Proceedings of the American Mathematical Society
\textbf{23} (1969), 679--688.

\bibitem{Karaliolios}
N.~Karaliolios,
\emph{Cohomological rigidity and the Anosov--Katok construction},
arXiv:1711.02732.
See the Introduction, p.~6, for the open question concerning
vanishing Lyapunov exponents.

\bibitem{Pesin}
Ya.~B.~Pesin,
\emph{Characteristic Lyapunov exponents and smooth ergodic theory},
Russian Mathematical Surveys
\textbf{32} (1977), no.~4, 55--114.

\end{thebibliography}

\end{document}
